\documentclass[11pt]{article}

% Plantilla basada en el formato de artículos de la ACL.
% Para la versión final, omite la opción "review".
\usepackage{acl}
\usepackage{times}
\usepackage{latexsym}
\usepackage[T1]{fontenc}
\usepackage[utf8]{inputenc}
\usepackage[spanish,es-nodecimaldot]{babel}
\usepackage{microtype}
\usepackage{inconsolata}
\usepackage{graphicx}
\usepackage{booktabs}
\usepackage{amsmath}

\title{Ampliación de la arquitectura de un modelo de lengua básico}

\author{Carlos V.\\
  Grado en Ingeniería en Inteligencia Artificial, Universitat d'Alacant\\
  \texttt{caviro50@gmail.com}}

\begin{document}
\maketitle

\begin{abstract}
Esta memoria describe la ampliación implementada sobre el modelo de lengua
base, el protocolo experimental empleado para validarla y los resultados
obtenidos. Sustituye este texto por un resumen de entre 100 y 150 palabras.
\end{abstract}

\section{Introducción}
El código de partida implementa los componentes de un transformer causal,
incluyendo la atención multi-cabeza, la normalización, la red \emph{feed-forward}
y la generación de texto. En este trabajo se amplía dicha arquitectura con
\textbf{[nombre de la técnica]}.

\section{Trabajo relacionado}
Describe brevemente la técnica elegida y los modelos de lengua que la emplean.
Incluye citas a las fuentes originales, por ejemplo \citep{raschka2024}.

\section{Metodología}
\subsection{Modelo base}
Documenta la configuración del modelo y los módulos que se modifican.

\subsection{Ampliación propuesta}
Explica la implementación, las decisiones de diseño y las diferencias respecto
al código base. Añade pseudocódigo o ecuaciones cuando resulten útiles.

\section{Evaluación experimental}
Especifica una validación reproducible: comparación numérica con una
implementación de referencia, medidas de memoria o velocidad, o entrenamiento
de un modelo pequeño. Indica los datos, las semillas, el hardware y las
métricas utilizados.

\begin{table}[t]
  \centering
  \begin{tabular}{lcc}
    \toprule
    Configuración & Métrica 1 & Métrica 2 \\
    \midrule
    Modelo base & -- & -- \\
    Modelo ampliado & -- & -- \\
    \bottomrule
  \end{tabular}
  \caption{Resultados de la evaluación. Sustituye los guiones por los valores obtenidos.}
  \label{tab:resultados}
\end{table}

\section{Resultados y discusión}
Presenta e interpreta los resultados de la Tabla~\ref{tab:resultados}. Indica
si las evidencias confirman el efecto esperado de la ampliación y comenta sus
limitaciones.

\section{Conclusiones}
Resume la contribución, los hallazgos principales y posibles mejoras futuras.

\bibliography{referencias}

\end{document}
